\section{Experiments}\label{sec:exp}

We ask three questions. Do the two backends behave as one engine wearing two geometries (E1, E4)? Does the shared classifier detect periodic structure, and how does it fail (E2)? Which mechanism actually stabilises the static layer (E3)?

\subsection{Experimental Setup}\label{subsec:setup}

\noindent\textbf{Harness and protocol.} All four experiments are driven by a standalone harness that compiles directly against the \code{strata\_core} sources, with no ROS~2, ament or colcon in the loop; it lives in \nolinkurl{paper/experiments/}. One command, \nolinkurl{bash paper/experiments/run_all.sh}, builds the core and its unit tests, checks the \code{ctest} suite (one target aggregating the 32 gtest cases, ten of them regression tests for the amplitude defect of Remark~\ref{rem:leak}), builds and runs the four executables, and writes the CSVs in \nolinkurl{paper/experiments/results/}; every table and plot in this paper is regenerated from those files (Appendix~\ref{app:repro}). Every stochastic experiment derives its \code{std::mt19937} generator from the constant seed 12345 plus a per-run offset recorded in code, and the seed is logged in each CSV. All magnitudes are therefore single-seed point estimates: we report no cross-seed dispersion, and gaps such as 0.9253 against 0.9758 illustrate a mechanism rather than a distributional estimate.

\noindent\textbf{What is compared.} There is no external baseline: none of the related systems (\S\ref{sec:related}) exposes a SLAM-free, ROS-free persistence core that could be driven by the same harness, so a single-number comparison would not be like for like. Instead E1 and E4 compare the two \sname\ backends, each driven in its native geometry by its own seeded stream (\code{grid2d} on the integer lattice, \code{voxel3d} continuously); because both call the identical engine, any divergence arises in $\kappa_b$, $\rho_b$ or their input sampling. E2 and E3 drive \code{LayeredMap} directly with no backend at all, which by \eqref{eq:objective} characterises the classifier for whichever backend supplies the ids. Each harness sets its own parameters in code (Table~\ref{tab:eparams}); in particular all four run with one window per tick ($L{=}1$), so \emph{window} and \emph{frame} coincide.

\noindent\textbf{Metrics.} Static-set precision, recall and F1 against a known wall set (E1, E3); Periodic-class true- and false-positive rates against ground truth and the dominant amplitude $a$ (E2); \emph{flicker}, the number of per-window toggles of a wall cell's Static flag (E3); mean latency per \code{integrate()} call, cell updates per second in \code{endWindow()}, final live-cell count and estimated memory at 56~B per cell (E4).

\noindent\textbf{Scenarios and hardware.} E1: a 161-cell static cross (the segments $y{=}60$ and $x{=}60$ over $[20,100]$) hit every window, plus 5, 25 or 100 fresh random movers per window, for 40 windows, on a $120{\times}120$ grid and a voxel map, both at resolution 1.0, once with periodicity off and once with it on at the shipped $T{=}24$, $H{=}2$. E2: three periodic doors (period 8 at 4-on/4-off and 2-on/6-off, period 4 at 2-on/2-off), a constant wall and four Bernoulli(0.5) aperiodic controls, with $T{=}8$, $H{=}3$, read out at 64 windows and, in a sweep, after every length $n=8$--$100$ (a fresh map per length). E3: 50 wall cells observed occupied with probability 0.6 per window among 10 movers per window, 80 windows, the same noise replayed for all 36 configurations of $\pgrad\in\{0.6,0.7,0.8,0.9\}\times\pdem\in\{0.3,0.5,\pgrad\}\times\lambda\in\{0.90,0.97,1.0\}$, with periodicity off and again with it on at the shipped defaults. All results are from the corrected engine; the same harness run against the v0.1.0 engine is kept for comparison (Table~\ref{tab:fix}). E4: 150 random hits per window for 20 windows at nominal extents $100^2$, $250^2$ and $500^2$, resolution 1.0, periodicity off. E4 alone reads the wall clock; its absolute microseconds come from one machine whose specification is not recorded in the committed artifacts and are not claimed to be portable, whereas its live-cell counts follow only from ray geometry and are deterministic.

\subsection{Main Results}\label{subsec:main}

% AUTO-GENERATED by tools/make_artifacts.py from paper/experiments/results/e1_static_quality.csv -- do not edit
\begin{table}[t]
\centering
\caption{Static-layer quality (E1) against the 161-cell ground-truth wall set. \emph{Recall~1 from} is the first window index $t$ at which recall reaches 1.0 (it never drops afterwards); precision, F1 and the predicted static-set size are at the final window $t{=}39$. Single seed (12345), so no dispersion is reported. Bold marks the better backend under high clutter.}
\label{tab:e1}
\small
\begin{tabular}{llcccc}
\toprule
Backend & Clutter (movers/window) & Recall 1 from & Precision & F1 & $|\hat{\mathcal{S}}|$ \\
\midrule
\texttt{grid2d} & low (5) & $t{=}2$ & 1.0000 & 1.0000 & 161 \\
\texttt{grid2d} & med (25) & $t{=}2$ & 0.9938 & 0.9969 & 162 \\
\texttt{grid2d} & high (100) & $t{=}2$ & 0.9253 & 0.9612 & 174 \\
\midrule
\texttt{voxel3d} & low (5) & $t{=}2$ & 1.0000 & 1.0000 & 161 \\
\texttt{voxel3d} & med (25) & $t{=}2$ & 0.9938 & 0.9969 & 162 \\
\texttt{voxel3d} & high (100) & $t{=}2$ & \textbf{0.9758} & \textbf{0.9877} & 165 \\
\bottomrule
\end{tabular}
\end{table}

\begin{figure}[t]
\centering
\begin{tikzpicture}
\begin{groupplot}[group style={group size=2 by 1, horizontal sep=14mm},
  width=0.47\linewidth, height=0.32\linewidth, xmin=0, xmax=39, ymin=0.93, ymax=1.005,
  xlabel={window $t$}, grid=major, grid style={gray!20}, tick label style={font=\scriptsize},
  label style={font=\small}, title style={font=\small},
  legend style={font=\scriptsize, at={(0.97,0.03)}, anchor=south east}, legend cell align=left,
  restrict y to domain=0.9:1.1, unbounded coords=discard]
\nextgroupplot[title={\code{grid2d}}, ylabel={static-set F1}]
\addplot[thick, blue!70!black] table[x=window,y=f1] {figures/data/e1_grid2d_low.dat};
\addplot[thick, orange!80!black, dashed] table[x=window,y=f1] {figures/data/e1_grid2d_med.dat};
\addplot[thick, red!70!black, densely dotted] table[x=window,y=f1] {figures/data/e1_grid2d_high.dat};
\legend{low (5/w), med (25/w), high (100/w)}
\nextgroupplot[title={\code{voxel3d}}]
\addplot[thick, blue!70!black] table[x=window,y=f1] {figures/data/e1_voxel3d_low.dat};
\addplot[thick, orange!80!black, dashed] table[x=window,y=f1] {figures/data/e1_voxel3d_med.dat};
\addplot[thick, red!70!black, densely dotted] table[x=window,y=f1] {figures/data/e1_voxel3d_high.dat};
\end{groupplot}
\end{tikzpicture}
\caption{\textbf{Static-layer F1 over time (E1).} F1 is 0 for $t<2$ (no cell has met $\Nmin{=}3$; not shown) and reaches 1.0 at $t{=}2$ for every backend and clutter level. Only the 100-movers-per-window regime accumulates false statics, more so on \code{grid2d}. Plotted directly from \code{e1\_static\_quality.csv}.}
\label{fig:e1}
\end{figure}


\noindent\textbf{Backend equivalence of the static layer (E1).} With $\Nmin{=}3$ a cell needs three touched windows before it may graduate, and recall jumps from 0 to 1.0 exactly at the third window ($t{=}2$) for all six backend--clutter combinations, then holds at 1.0 through $t{=}39$: clutter never demotes a wall cell (Table~\ref{tab:e1}, Figure~\ref{fig:e1}). Table~\ref{tab:e1} is the periodicity-off run. With periodicity on at the shipped $T{=}24$, $H{=}2$ the run is identical window for window: the wall is occupied whenever it is touched, so its centred amplitude is 0 (Remark~\ref{rem:leak}). The v0.1.0 engine instead lost the whole wall from the Static layer in windows 29--40 ($t=28$--$39$), in all six runs, and ended with Static recall 0 (Table~\ref{tab:fix}). The only errors are false Static cells created when a random mover re-hits one cell often enough to graduate on its own. At 5 movers per window this never happens and precision stays at 1.0; at 25 it saturates at one spurious cell (precision 0.9938); at 100 false statics keep accumulating, from the first extra cell at $t{=}17$ (\code{grid2d}) or $t{=}25$ (\code{voxel3d}) to 174 or 165 predicted cells at $t{=}39$. This yields the one material gap between backends, precision 0.9253 against 0.9758. Both runs apply the identical graduation rule, each to its own natively sampled mover stream (movers on the integer lattice for \code{grid2d}, continuous for \code{voxel3d}, both confined to one $z$-plane), so the gap reflects quantisation and sampling rather than the classifier; we do not claim it isolates a single geometric mechanism.

% AUTO-GENERATED by tools/make_artifacts.py from paper/experiments/results/e2_classification.csv, e2_summary.csv -- do not edit
\begin{table}[t]
\centering
\caption{Periodicity detection (E2) after 64 windows, $a_{\min}{=}0.3$, $T{=}8$, $H{=}3$. \emph{Ref.\ amplitude} is a non-pruning reference model fed the identical occupancy stream. Periodic TPR $=2/3$, FPR $=1/5$.}
\label{tab:e2}
\small
\begin{tabular}{lcccc}
\toprule
Probe cell & Ground truth & Final class & Ref.\ amplitude & Outcome \\
\midrule
\texttt{door\_p8\_4on4off} (50\% duty) & periodic & Periodic & 0.653 & correct \\
\texttt{door\_p8\_2on6off} (25\% duty) & periodic & Transient & 0.462 & \textbf{miss} \\
\texttt{door\_p4\_2on2off} (50\% duty) & periodic & Periodic & 0.707 & correct \\
\texttt{wall\_constant} & non-periodic & Static & $\approx 0$ & correct \\
\texttt{aperiodic\_0} & non-periodic & Transient & 0.159 & correct \\
\texttt{aperiodic\_1} & non-periodic & Periodic & 0.329 & \textbf{false pos.} \\
\texttt{aperiodic\_2} & non-periodic & Transient & 0.139 & correct \\
\texttt{aperiodic\_3} & non-periodic & Transient & 0.251 & correct \\
\bottomrule
\end{tabular}
\end{table}

\begin{figure}[t]
\centering
\begin{tikzpicture}
\begin{axis}[width=0.82\linewidth, height=0.40\linewidth, xmode=log, log basis x=2,
  xtick={6,8,12,16,24,32,48,64}, xticklabels={6,8,12,16,24,32,48,64},
  xmin=5.5, xmax=70, ymin=-0.02, ymax=0.92,
  xlabel={observation length $n$ (touched windows)}, ylabel={dominant amplitude $a$},
  grid=major, grid style={gray!20}, tick label style={font=\scriptsize}, label style={font=\small},
  legend style={font=\scriptsize, at={(1.02,0.5)}, anchor=west}, legend cell align=left]
\addplot[thick, blue!70!black, mark=*, mark size=1.2pt] table[x=len,y=amp] {figures/data/e2_door_p8_4on4off.dat};
\addplot[thick, blue!40, mark=square*, mark size=1.2pt] table[x=len,y=amp] {figures/data/e2_door_p4_2on2off.dat};
\addplot[thick, teal!70!black, dashed, mark=triangle*, mark size=1.4pt] table[x=len,y=amp] {figures/data/e2_door_p8_2on6off.dat};
\addplot[thick, black, mark=x] table[x=len,y=amp] {figures/data/e2_wall_constant.dat};
\addplot[red!70!black, thick, densely dotted, mark=o, mark size=1.2pt] table[x=len,y=amp] {figures/data/e2_aperiodic_1.dat};
\addplot[gray, thin, mark=none] table[x=len,y=amp] {figures/data/e2_aperiodic_0.dat};
\addplot[gray, thin, mark=none, forget plot] table[x=len,y=amp] {figures/data/e2_aperiodic_2.dat};
\addplot[gray, thin, mark=none, forget plot] table[x=len,y=amp] {figures/data/e2_aperiodic_3.dat};
\addplot[red, thick, dashed, domain=5.5:70, samples=2, forget plot] {0.3};
\addplot[red!60!black, thick, dotted, forget plot] table[x=len,y=astar] {figures/data/e2_threshold.dat};
\node[font=\scriptsize, anchor=south west, text=red!60!black] at (axis cs:9,0.80) {$a^\star(n)$};
\draw[gray!70, dashed] (axis cs:8,-0.02) -- (axis cs:8,0.92);
\node[font=\scriptsize, anchor=south west, text=red] at (axis cs:40,0.30) {$\amin{=}0.3$};
\node[font=\scriptsize, anchor=north west, text=gray] at (axis cs:8,0.9) {$n{=}T$};
\legend{door p8 4/4, door p4 2/2, door p8 2/6, wall, aperiodic 1, aperiodic 0/2/3}
\end{axis}
\end{tikzpicture}
\caption{\textbf{FreMEn-lite amplitude against observation length (E2)}, from a non-pruning reference model of the corrected engine fed each probe cell's occupancy stream ($T{=}8$, $H{=}3$). The $n\ge T$ gate forces $a{=}0$ at $n{=}6$; both 50\%-duty doors settle well above $\amin$, and the 25\%-duty door would too, yet the live pipeline prunes it first (Proposition~\ref{prop:race}). The constant wall is exactly 0 at every length because the coefficients are centred (Remark~\ref{rem:leak}); the Bernoulli controls hover near $\amin$. The dotted curve $a^\star(n)$ is the amplitude the calibrated test of Proposition~\ref{prop:chernoff} requires at $\delta{=}0.1$, $H{=}3$ under uniform phase coverage (Eq.~\eqref{eq:astar}); a cell is Periodic only above both curves. Plotted from \code{e2\_amplitude\_vs\_length.csv}.}
\label{fig:e2}
\end{figure}

\begin{figure}[t]
\centering
\begin{tikzpicture}
\begin{axis}[name=fpr, width=0.82\linewidth, height=0.32\linewidth,
  xmin=7, xmax=101, ymin=-0.03, ymax=0.85, xtick={8,16,24,32,40,48,56,64,72,80,88,96},
  xticklabels={}, ylabel={Periodic FPR (of 5)},
  grid=major, grid style={gray!20}, tick label style={font=\scriptsize}, label style={font=\small},
  legend style={font=\scriptsize, at={(0.5,1.03)}, anchor=south, legend columns=3}, legend cell align=left]
\addplot[const plot, thick, gray, densely dashed] table[x=len,y=fpr] {figures/data/e2_fpr_pre.dat};
\addplot[const plot, thick, orange!80!black] table[x=len,y=fpr] {figures/data/e2_fpr_mid.dat};
\addplot[const plot, thick, violet!70!black] table[x=len,y=fpr] {figures/data/e2_fpr_one.dat};
\addplot[const plot, very thick, blue!70!black] table[x=len,y=fpr] {figures/data/e2_fpr_post.dat};
\addplot[const plot, thin, violet!70!black, densely dotted] table[x=len,y=reffpr] {figures/data/e2_fpr_one.dat};
\addplot[const plot, thin, blue!70!black, densely dotted] table[x=len,y=reffpr] {figures/data/e2_fpr_post.dat};
\draw[gray!70, dashed] (axis cs:64,-0.03) -- (axis cs:64,0.85);
\legend{v0.1.0, centred, single read-out, spent, single (no pruning), spent (no pruning)}
\end{axis}
\begin{axis}[at={(fpr.south west)}, anchor=north west, yshift=-4pt,
  width=0.82\linewidth, height=0.24\linewidth,
  xmin=7, xmax=101, ymin=-0.05, ymax=1.05, xtick={8,16,24,32,40,48,56,64,72,80,88,96},
  ytick={0,0.333,0.667,1}, yticklabels={0,1/3,2/3,1},
  xlabel={read-out length $n$ (windows)}, ylabel={TPR (of 3)},
  grid=major, grid style={gray!20}, tick label style={font=\scriptsize}, label style={font=\small}]
\addplot[const plot, thick, gray, densely dashed] table[x=len,y=tpr] {figures/data/e2_fpr_pre.dat};
\addplot[const plot, thick, orange!80!black] table[x=len,y=tpr] {figures/data/e2_fpr_mid.dat};
\addplot[const plot, thick, violet!70!black] table[x=len,y=tpr] {figures/data/e2_fpr_one.dat};
\addplot[const plot, very thick, blue!70!black] table[x=len,y=tpr] {figures/data/e2_fpr_post.dat};
\addplot[const plot, thin, blue!70!black, densely dotted] table[x=len,y=reftpr] {figures/data/e2_fpr_post.dat};
\draw[gray!70, dashed] (axis cs:64,-0.05) -- (axis cs:64,1.05);
\end{axis}
\end{tikzpicture}
\caption{\textbf{Periodic false- and true-positive rates against read-out length (E2)}, from the real \code{LayeredMap}, a fresh map per length, $T{=}8$, $H{=}3$, $\amin{=}0.3$, for the four engine versions of Table~\ref{tab:fix}. The dotted lines apply the single-read-out and the spent predicate to a non-pruning reference model fed the same streams (lower panel: spent only). The dashed vertical marks the $n{=}64$ read-out of Table~\ref{tab:e2}. The v0.1.0 and centred TPR curves coincide at 2/3. Plotted from \code{e2\_rates\_vs\_length.csv} and its three archived counterparts.}
\label{fig:e2fpr}
\end{figure}

% AUTO-GENERATED by tools/make_artifacts.py from paper/experiments/results/e1_static_quality{,_periodic}.csv, e2_summary.csv, e2_rates_vs_length.csv, e3_sensitivity{,_periodic}.csv, pre_fix_2026-09-28/, pre_calibration_2026-09-28/ and pre_spending_2026-09-28/ -- do not edit
\begin{table}[t]
\centering
\caption{Effect of the three corrections to the periodicity test, same harness, seeds and parameters. \emph{v0.1.0}: uncentred amplitude and the amplitude-only rule $a\ge\amin$ (results in \nolinkurl{results/pre\_fix\_2026-09-28/}); \emph{Centred}: centred amplitude, amplitude-only rule (\nolinkurl{results/pre\_calibration\_2026-09-28/}); \emph{Single}: the significance test of Proposition~\ref{prop:chernoff} at $\delta{=}0.1$ at every read-out (\nolinkurl{results/pre\_spending\_2026-09-28/}); \emph{Spent}: the level $\delta{=}0.2$ spent over the touch count (Proposition~\ref{prop:spend}, current engine). E1 and E3 with periodicity on use the shipped $T{=}24$, $H{=}2$, $\amin{=}0.3$; E2 uses $T{=}8$, $H{=}3$. E1 counts are identical in all six backend--clutter runs.}
\label{tab:fix}
\footnotesize
\setlength{\tabcolsep}{4pt}
\begin{tabular}{lccccc}
\toprule
& Periodicity & \multicolumn{4}{c}{Periodicity on} \\
\cmidrule(lr){3-6}
Measure & off & v0.1.0 & Centred & Single & Spent \\
\midrule
E1: windows $t\ge3$ with Static recall $<1$ & 0 windows & 12 ($t{=}28$--$39$) & 0 windows & 0 windows & 0 windows \\
E1: Static recall at $t{=}39$ & 1 & 0 & 1 & 1 & 1 \\
\midrule
E2 at $n{=}64$: TPR / FPR & -- & 2/3 / 1/5 & 2/3 / 2/5 & 2/3 / 0/5 & 2/3 / 0/5 \\
E2, $n=8$--$100$: mean FPR & -- & 0.363 & 0.331 & 0.000 & 0.000 \\
E2, $n=8$--$100$: FP count range (of 5) & -- & 0--4 & 0--4 & 0--0 & 0--0 \\
E2, $n=8$--$100$: lengths with no FP & -- & 10 of 93 & 7 of 93 & 93 of 93 & 93 of 93 \\
E2, $n=8$--$100$: lengths with constant wall Periodic & -- & 4 of 93 & 0 of 93 & 0 of 93 & 0 of 93 \\
E2, $n=8$--$100$: mean TPR & -- & 0.667 & 0.667 & 0.624 & 0.563 \\
E2: TPR 2/3 at every $n\ge$ & -- & 8 & 8 & 15 & 25 \\
E2, $n=8$--$100$: lengths with FP, no pruning & -- & -- & -- & 25 of 93 & 0 of 93 \\
\midrule
E3: total flicker, 36 configurations & 3672 & 7132 & 5472 & 3792 & 3672 \\
E3: flicker at $\pgrad{=}0.9$, $\pdem{=}0.3$, $\lambda{=}0.90$ & 54 & 138 & 102 & 58 & 54 \\
E3: flicker at $\pgrad{=}\pdem{=}0.9$, $\lambda{=}0.90$ & 574 & 582 & 596 & 576 & 574 \\
\bottomrule
\end{tabular}
\end{table}


\noindent\textbf{Periodicity detection (E2).} Both 50\%-duty doors are classified Periodic with amplitudes 0.653 and 0.707, well above $\amin{=}0.3$ (Table~\ref{tab:e2}). The constant wall is Static with amplitude exactly 0 at every length, and two of the four aperiodic controls end Transient. At the $n{=}64$ read-out the Periodic class therefore reaches a true-positive rate of 2/3 and a false-positive rate of 2/5. The v0.1.0 engine scored 1/5 at this length, so the corrected engine is \emph{worse} at the read-out previously reported, and we do not treat either value as representative. Figure~\ref{fig:e2fpr} reads the pipeline out at every length from 8 to 100. The true-positive rate is 2/3 at every length both before and after the fix, but the false-positive count ranges over 0--4 of 5 in both versions, with a mean false-positive rate of 0.363 before the fix and 0.331 after. The fix removes the constant wall from that count (it was Periodic at 4 of 93 lengths, now none). The remaining false positives are Bernoulli controls, and centring cannot suppress them because they genuinely vary. Two effects combine. First, $\amin{=}0.3$ is not calibrated against sampling noise: for a Bernoulli(0.5) cell observed $n$ times each centred coefficient has standard deviation $\sqrt{0.5/n}$, which is 0.25 at $n{=}T{=}8$. Second, pruning resets $\boldsymbol{\phi}_i$, so the live pipeline often estimates the amplitude from only the touches since a cell's last re-creation. At $n{=}64$, \code{aperiodic\_0} has reference amplitude 0.159 over all 64 windows but is Periodic in the pipeline, while \code{aperiodic\_1} has reference amplitude 0.329, 0.029 above $\amin$. A significance-calibrated threshold is future work (\S\ref{sec:disc}). The amplitude curves confirm the gate of \eqref{eq:periodic}: for the 4-on/4-off door the amplitude is exactly 0 at $n{=}6<T$, 0.653 at every sampled multiple of $T$ ($n=8,16,\dots,64$) and 0.581 at $n{=}12$, where the uncentred v0.1.0 amplitude had overshot to 0.871 (Figure~\ref{fig:e2}). The miss, the 25\%-duty door, is not a modelling failure: a non-pruning reference model fed the same stream reaches amplitude 0.462, above $\amin$. The live pipeline never gets there, for a reason that follows from the update rule.

\begin{proposition}[Prune/maturity race]\label{prop:race}
Under the E2 parameters ($\lhit{=}1$, $\lmiss{=}{-}1$, $\lambda{=}1$, $\pprune{=}0.05$, $T{=}8$), a cell touched every window with the 2-on/6-off pattern is erased by \eqref{eq:prune} within at most 7 consecutive touched windows of its creation. Hence $n_i<T$ throughout every lifetime, $a_i=0$, and the cell is never classified Periodic.
\end{proposition}
\begin{proof}
With $\lambda{=}1$ and unit increments, $\ell_i$ is an integer walk from $0$ (the clamp at $\pm5$ never binds below). Since $\sigma(-2)\approx0.119\ge\pprune>\sigma(-3)\approx0.047$, erasure happens exactly when $\ell_i$ first reaches $-3$, provided $a_i<\amin$, which holds while $n_i<T$. If the cell is created with $j\ge3$ vacant windows left in its cycle, it reaches $-3$ after 3 windows. If $j\le2$ vacant windows remain, $\ell_i$ reaches $-j$, then $2-j$ after the two occupied windows, then $-3$ after $5-j$ further vacant windows: $7$ windows in total. If it is created in an occupied window, $\ell_i\le2$ after at most two hits and $-3$ after five more misses, again at most 7. In every case the lifetime is at most $7<T{=}8$ touched windows, so $a_i=0$ by \eqref{eq:periodic} and $\pi_i$ is false. Erasure also needs $\lnot g_i$, which holds because $\ell_i\le2$ throughout, so $p_i\le\sigma(2)\approx0.881<\pgrad{=}0.9$ and the cell never graduates. Erasure discards $\boldsymbol{\phi}_i$, so each re-created cell starts again from $n_i=0$.
\end{proof}
The race couples the prune rule \eqref{eq:prune} to the touched-window gate; the reference model, which never prunes, is not subject to it.

\noindent\textbf{Per-dimension cost of the shared engine (E4).} At every tested extent \code{grid2d}'s \code{integrate()} is cheaper than \code{voxel3d}'s at the same nominal extent and hit rate, by 15.0$\times$ at $100^2$, 10.8$\times$ at $250^2$ and 5.5$\times$ at $500^2$ (Table~\ref{tab:e4}, Figure~\ref{fig:e4}). The live-cell counts explain the gap: \code{voxel3d} ends each run holding 10.0$\times$, 6.5$\times$ and 4.4$\times$ as many cells, because its half-voxel ray march spawns a hash entry for every traversed free voxel, whereas exact 2D Bresenham clears a bounded strip. Both ratios narrow together as extent grows and rays cover a smaller share of the volume. Memory is exactly what a flat 56~B per cell predicts, tracking live-cell count with no additional constant, and reaches 37.0~MB for the largest 3D configuration (661,410 cells). The cost difference is geometry, not the engine.

% AUTO-GENERATED by tools/make_artifacts.py from paper/experiments/results/e4_throughput.csv -- do not edit
\begin{table}[t]
\centering
\caption{Cost of the shared engine under each backend (E4): 150 random hits per window, 20 windows, resolution 1.0, periodicity off. Latency and throughput are single-machine wall-clock samples; live-cell counts, ratios and memory (at 56~B per cell) are geometry-determined and reproducible. Ratios are voxel3d over grid2d at equal extent; bold marks the cheaper backend.}
\label{tab:e4}
\small
\begin{tabular}{llrrrr}
\toprule
Backend & Extent & $\mu$s / \texttt{integrate} $\downarrow$ & Cell updates/s & Live cells & Memory \\
\midrule
\texttt{grid2d} & $100^2$ & \textbf{320.8} & $9.41{\times}10^{6}$ & 8{,}553 & 479~KB \\
\texttt{voxel3d} & $100^2$ & 4{,}796.8 & $7.06{\times}10^{6}$ & 85{,}399 & 4.78~MB \\
\multicolumn{2}{l}{\quad\emph{ratio}} & \emph{15.0$\times$} & & \emph{10.0$\times$} & \\
\midrule
\texttt{grid2d} & $250^2$ & \textbf{1{,}702.6} & $1.07{\times}10^{7}$ & 45{,}657 & 2.56~MB \\
\texttt{voxel3d} & $250^2$ & 18{,}469.5 & $4.46{\times}10^{6}$ & 295{,}300 & 16.5~MB \\
\multicolumn{2}{l}{\quad\emph{ratio}} & \emph{10.8$\times$} & & \emph{6.5$\times$} & \\
\midrule
\texttt{grid2d} & $500^2$ & \textbf{6{,}548.9} & $8.64{\times}10^{6}$ & 150{,}096 & 8.41~MB \\
\texttt{voxel3d} & $500^2$ & 35{,}958.3 & $4.45{\times}10^{6}$ & 661{,}410 & 37.0~MB \\
\multicolumn{2}{l}{\quad\emph{ratio}} & \emph{5.5$\times$} & & \emph{4.4$\times$} & \\
\bottomrule
\end{tabular}
\end{table}

\begin{figure}[t]
\centering
\begin{tikzpicture}
\begin{groupplot}[group style={group size=2 by 1, horizontal sep=18mm},
  width=0.46\linewidth, height=0.34\linewidth, ybar=1pt, /pgf/bar width=7pt, ymode=log,
  symbolic x coords={100,250,500}, xtick=data, xticklabels={$100^2$,$250^2$,$500^2$},
  xlabel={nominal extent}, enlarge x limits=0.25, grid=major, grid style={gray!20},
  tick label style={font=\scriptsize}, label style={font=\small},
  legend style={font=\scriptsize, at={(0.03,0.97)}, anchor=north west}, legend cell align=left]
\nextgroupplot[ylabel={$\mu$s per \code{integrate()}}]
\addplot[fill=blue!55, draw=blue!70!black] table[x=extent, y=us, discard if not={backend}{grid2d}] {figures/data/e4.dat};
\addplot[fill=orange!60, draw=orange!80!black] table[x=extent, y=us, discard if not={backend}{voxel3d}] {figures/data/e4.dat};
\legend{\code{grid2d}, \code{voxel3d}}
\nextgroupplot[ylabel={live cells at end of run}]
\addplot[fill=blue!55, draw=blue!70!black] table[x=extent, y=cells, discard if not={backend}{grid2d}] {figures/data/e4.dat};
\addplot[fill=orange!60, draw=orange!80!black] table[x=extent, y=cells, discard if not={backend}{voxel3d}] {figures/data/e4.dat};
\end{groupplot}
\end{tikzpicture}
\caption{\textbf{Per-call cost and live-cell count (E4)}, log scale. The latency gap between backends narrows with extent in step with the live-cell gap, which is set by the 3D free-space ray march rather than by the shared engine. Plotted from \code{e4\_throughput.csv}.}
\label{fig:e4}
\end{figure}


\subsection{Ablation Studies}\label{subsec:ablation}

% AUTO-GENERATED by tools/make_artifacts.py from paper/experiments/results/e3_sensitivity.csv -- do not edit
\begin{table}[t]
\centering
\caption{Hysteresis and decay ablation (E3), selected rows of the 36-configuration sweep on identical replayed noise (full sweep in Table~\ref{tab:e3full}). Band $=p_{\mathrm{grad}}-p_{\mathrm{dem}}$; flicker counts per-window \texttt{isStatic} toggles over 80 windows. Bold marks the lowest flicker in the sweep and the worst configuration.}
\label{tab:e3}
\small
\begin{tabular}{cccccc}
\toprule
$p_{\mathrm{grad}}$ & $p_{\mathrm{dem}}$ & Band & $\lambda$ & Flicker $\downarrow$ & F1 $\uparrow$ \\
\midrule
0.9 & 0.9 (deg.) & 0.0 & 0.90 & \textbf{574} & \textbf{0.810} \\
0.8 & 0.8 (deg.) & 0.0 & 0.90 & 343 & 0.947 \\
0.7 & 0.7 (deg.) & 0.0 & 0.90 & 250 & 0.980 \\
\midrule
0.9 & 0.9 (deg.) & 0.0 & 0.97 & 150 & 1.000 \\
0.9 & 0.9 (deg.) & 0.0 & 1.00 & 122 & 1.000 \\
\midrule
0.9 & 0.3 & 0.6 & 0.90 & 54 & 1.000 \\
0.9 & 0.3 & 0.6 & 0.97 & \textbf{52} & \textbf{1.000} \\
0.9 & 0.3 & 0.6 & 1.00 & \textbf{52} & \textbf{1.000} \\
\bottomrule
\end{tabular}
\end{table}

\begin{figure}[t]
\centering
\begin{tikzpicture}
\begin{groupplot}[group style={group size=3 by 1, horizontal sep=4mm, yticklabels at=edge left},
  width=0.36\linewidth, height=0.36\linewidth, enlargelimits=false,
  colormap={flick}{rgb255=(255,255,255) rgb255=(253,208,162) rgb255=(230,85,13) rgb255=(127,39,4)},
  point meta min={1.716}, point meta max={2.759},
  xtick={0,1,2}, xticklabels={0.3,0.5,$=\!\pgrad$}, ytick={0,1,2,3}, yticklabels={0.6,0.7,0.8,0.9},
  xlabel={$\pdem$}, tick label style={font=\scriptsize}, label style={font=\small}, title style={font=\small},
  nodes near coords style={anchor=center}]
\nextgroupplot[title={$\lambda{=}0.90$}, ylabel={$\pgrad$}]
\addplot[matrix plot*, mesh/cols=3, point meta={log10(\thisrow{flicker})},
  visualization depends on={\thisrow{flicker} \as \fl}, nodes near coords={\scriptsize\pgfmathprintnumber[fixed,precision=0]{\fl}}] table[x=col,y=row] {figures/data/e3_decay_0.9.dat};
\nextgroupplot[title={$\lambda{=}0.97$}]
\addplot[matrix plot*, mesh/cols=3, point meta={log10(\thisrow{flicker})},
  visualization depends on={\thisrow{flicker} \as \fl}, nodes near coords={\scriptsize\pgfmathprintnumber[fixed,precision=0]{\fl}}] table[x=col,y=row] {figures/data/e3_decay_0.97.dat};
\nextgroupplot[title={$\lambda{=}1.0$}]
\addplot[matrix plot*, mesh/cols=3, point meta={log10(\thisrow{flicker})},
  visualization depends on={\thisrow{flicker} \as \fl}, nodes near coords={\scriptsize\pgfmathprintnumber[fixed,precision=0]{\fl}}] table[x=col,y=row] {figures/data/e3_decay_1.dat};
\end{groupplot}
\end{tikzpicture}
\caption{\textbf{Flicker transitions across the E3 sweep} (colour: $\log_{10}$ flicker; annotation: raw count). The zero-band column ($\pdem{=}\pgrad$) holds the worst cell of every panel; widening the band suppresses flicker far more than slowing decay does. Plotted from \code{e3\_sensitivity.csv}.}
\label{fig:e3}
\end{figure}


\noindent\textbf{Hysteresis band versus decay (E3).} E3 ablates the two stabilisers of the static layer on identical replayed noise (Table~\ref{tab:e3}, full sweep in Table~\ref{tab:e3full}). Removing the hysteresis band has the largest effect in the sweep: at $\pgrad{=}\pdem{=}0.9$ and $\lambda{=}0.90$, the worst configuration, the wall layer flickers 574 times and final F1 drops to 0.810 (recall 0.68), whereas the same noise and decay with the band widened to 0.6 ($\pdem{=}0.3$) flickers 54 times at F1 1.000, a more than tenfold reduction attributable to the band alone. Among zero-band configurations at $\lambda{=}0.90$, flicker rises with the coincident threshold (250 at 0.7, 343 at 0.8, 574 at 0.9), even though a higher threshold is usually the more conservative choice: a threshold pair with no gap is not a substitute for a band wherever it sits. Decay is second order. Holding the worst pair fixed, flicker falls as forgetting slows (574 at $\lambda{=}0.90$, 150 at 0.97, 122 at 1.0), because faster decay lets noisy log-odds cross the coincident threshold more often; with the widest band, decay barely matters (52--54 toggles across all three settings). Hysteresis, not decay, is what keeps the static layer stable (Figure~\ref{fig:e3}).

\noindent\textbf{Periodicity on (E3).} Rerunning the sweep with periodicity on at the shipped defaults leaves final F1 unchanged in every configuration, but raises flicker in every one (last column of Table~\ref{tab:e3full}). The total over the 36 configurations is 5472, against 3672 with periodicity off and 7132 with the v0.1.0 engine; the widest-band configuration at $\lambda{=}0.90$ goes from 54 to 102 toggles (138 before the fix). These walls are free in 40\% of their windows by construction, so their centred amplitude is sampling noise that crosses $\amin$ now and then, and the $\lor\,\pi_i$ demotion, which has no band of its own (\S\ref{subsec:schmitt}), then toggles the wall out of the Static layer. The hysteresis band protects the occupancy path only. A persistence rule on $\pi_i$ or a noise-calibrated $\amin$ would be needed to extend that protection to the periodicity path, and we have not evaluated either.
